% © 2026 Ing. David Jaroš — CC BY-NC-ND 4.0
%
% This work is licensed under a Creative Commons Attribution-NonCommercial-NoDerivatives
% 4.0 International License (CC BY-NC-ND 4.0).
%
% License History: Earlier drafts (up to v0.3) were released under CC BY 4.0.
% From v0.4 onward, all material is released under CC BY-NC-ND 4.0 to protect
% the integrity of the theoretical work during ongoing academic development.
%
% See LICENSE.md for full license text.

% SU(3) Gauge Connection — Three-Qubit Encoding
% Version: 1.0
% Date: 2026-04-27
% Status: Canonical — equivalence proof at representation level
% Author: Ing. David Jaroš

\ifdefined\INCLUDEMODE
  \def\UBTSUThreeEndDocument{}
\else
  \documentclass[12pt,a4paper]{article}
  \usepackage[a4paper, margin=2.5cm]{geometry}
  \usepackage{amsmath, amssymb, amsthm}
  \usepackage{mathtools}
  \usepackage{hyperref}
  \usepackage{xcolor}
  \usepackage{tcolorbox}
  \tcbuselibrary{skins,breakable}
  \usepackage{booktabs}
  \usepackage{array}

  \newtheorem{theorem}{Theorem}[section]
  \newtheorem{lemma}[theorem]{Lemma}
  \newtheorem{proposition}[theorem]{Proposition}
  \newtheorem{corollary}[theorem]{Corollary}
  \theoremstyle{definition}
  \newtheorem{definition}[theorem]{Definition}
  \theoremstyle{remark}
  \newtheorem{remark}[theorem]{Remark}

  \newcommand{\C}{\mathbb{C}}
  \newcommand{\HH}{\mathbb{H}}
  \newcommand{\Hcolor}{\mathcal{H}_{\mathrm{color}}}
  \newcommand{\Hthreeq}{\mathcal{H}_{3q}}
  \newcommand{\Pcolor}{P_{\mathrm{color}}}

  \title{\textbf{SU(3) Gauge Connection: Three-Qubit Encoding}\\
         \large Equivalence at the Level of Lie Algebra Representation
         and Covariant Derivative}
  \author{Ing.\ David Jaro\v{s}}
  \date{2026-04-27}

  \begin{document}
  \def\UBTSUThreeEndDocument{\end{document}}
  \maketitle

  \begin{abstract}
  We establish a precise equivalence between the standard UBT SU(3) gauge
  connection $A_\mu = A_\mu^a \lambda_a/2$ and a three-qubit encoded gauge
  connection $A_\mu^{3q} = A_\mu^a G_a/2$, where $G_a$ are the Gell--Mann
  matrices lifted to the one-hot (color) subspace of
  $\Hthreeq = \C^2 \otimes \C^2 \otimes \C^2$.
  The equivalence is representation-level: both descriptions carry the same
  $\mathfrak{su}(3)$ Lie algebra and the same local gauge covariant derivative on
  the color subspace $\Hcolor = \mathrm{span}\{|100\rangle, |010\rangle,
  |001\rangle\}$.  The full three-qubit space contains an orthogonal leakage
  sector $\Hcolor^\perp$, which must be constrained or interpreted separately.
  \end{abstract}
\fi

%──────────────────────────────────────────────────────────────────────────────
\section{Setup and Notation}
\label{sec:su3qubit:setup}
%──────────────────────────────────────────────────────────────────────────────

\begin{tcolorbox}[colback=blue!5!white,colframe=blue!40!black,
  title=Status and scope]
\textbf{Scope}: This file proves equivalence at the Lie algebra representation
and covariant derivative level.  It does \emph{not} claim physical derivation of
full QCD dynamics, confinement, anomaly cancellation, or the fermion
representation content.

\medskip
\noindent\textbf{Canonical SU(3) source}: Theorems G.A--G.D in
\texttt{canonical/interactions/sm\_gauge.tex}; the involution derivation
in \texttt{canonical/su3\_derivation/su3\_from\_involutions.tex}.

\medskip
\noindent\textbf{Three-qubit source}: Geometric analysis in
\texttt{research\_tracks/triqubit\_su3\_geometry.md};
one-hot embedding in
\texttt{research\_tracks/THEORY\_COMPARISONS/su3\_qubit\_mapping/su3\_qubit\_core/mapping.py}.
\end{tcolorbox}

\paragraph{Conventions.}
We follow the standard normalisation $\mathrm{Tr}[\lambda_a \lambda_b] =
2\delta_{ab}$ for Gell-Mann matrices.  Indices $a,b,c \in \{1,\ldots,8\}$ label
adjoint generators; $i,j \in \{1,2,3\}$ label the fundamental color triplet.
The strong coupling is $g$ (or $g_s$ when disambiguation is needed).

%──────────────────────────────────────────────────────────────────────────────
\section{Definitions}
\label{sec:su3qubit:defs}
%──────────────────────────────────────────────────────────────────────────────

\begin{definition}[Three-qubit Hilbert space]
\label{def:H3q}
The \emph{three-qubit Hilbert space} is
\[
  \Hthreeq = \C^2 \otimes \C^2 \otimes \C^2 \cong \C^8,
\]
with computational basis $\{|q_1 q_2 q_3\rangle \mid q_i \in \{0,1\}\}$.
\end{definition}

\begin{definition}[Color subspace]
\label{def:Hcolor}
The \emph{color subspace} (one-hot sector) is
\[
  \Hcolor = \mathrm{span}\{|100\rangle,\, |010\rangle,\, |001\rangle\}
  \subset \Hthreeq.
\]
The three basis states are identified with the three quark color charges:
\begin{align*}
  |100\rangle &\longleftrightarrow |r\rangle \quad \text{(red)}, \\
  |010\rangle &\longleftrightarrow |g\rangle \quad \text{(green)}, \\
  |001\rangle &\longleftrightarrow |b\rangle \quad \text{(blue)}.
\end{align*}
\end{definition}

\begin{definition}[Color projector]
\label{def:Pcolor}
The orthogonal projector onto $\Hcolor$ is
\[
  \Pcolor = |100\rangle\langle 100| + |010\rangle\langle 010|
            + |001\rangle\langle 001| \;\in\; \mathrm{End}(\C^8).
\]
\end{definition}

\begin{definition}[Encoded generators]
\label{def:Ga}
For each Gell-Mann matrix $\lambda_a$ ($a = 1,\ldots,8$) acting on $\C^3$, the
\emph{encoded generator} $G_a \in \mathrm{End}(\C^8)$ is defined by
\[
  G_a = P\, \lambda_a\, P^\dagger,
\]
where $P: \C^3 \to \C^8$ is the isometric one-hot embedding
\[
  P = \begin{pmatrix}0 & 0 & 0 \\ 0 & 0 & 1 \\ 0 & 1 & 0 \\ 0 & 0 & 0 \\
                     1 & 0 & 0 \\ 0 & 0 & 0 \\ 0 & 0 & 0 \\ 0 & 0 & 0\end{pmatrix}
\]
mapping $|r\rangle \mapsto |100\rangle$, $|g\rangle \mapsto |010\rangle$,
$|b\rangle \mapsto |001\rangle$.

Equivalently, $G_a$ acts as $\lambda_a$ on $\Hcolor$ and as zero on
$\Hcolor^\perp$:
\[
  G_a\big|_{\Hcolor} = \lambda_a, \qquad
  G_a\big|_{\Hcolor^\perp} = 0.
\]
\end{definition}

\begin{remark}[Generator normalisation]
The matrices $G_a$ lift the full Gell--Mann matrices, not the conventionally
normalised Lie generators $\lambda_a/2$.  Consequently
$[G_a,G_b]=2if_{abc}G_c$, while gauge connections and exponentials use
$G_a/2$.  This distinction is kept explicit below.
\end{remark}

\begin{definition}[Encoded gauge connection and covariant derivative]
\label{def:Amu3q}
Given real gauge functions $A_\mu^a(x)$, define
\[
  A_\mu^{3q}(x) = \sum_{a=1}^{8} A_\mu^a(x)\, \frac{G_a}{2}
  \;\in\; \mathrm{End}(\C^8),
\]
and the \emph{encoded covariant derivative}
\[
  D_\mu^{3q} = \partial_\mu + ig\, A_\mu^{3q}.
\]
\end{definition}

%──────────────────────────────────────────────────────────────────────────────
\section{Basis and Generator Mapping}
\label{sec:su3qubit:tables}
%──────────────────────────────────────────────────────────────────────────────

\subsection{Color Basis Mapping}

\begin{table}[h]
\centering
\begin{tabular}{cccc}
\toprule
\textbf{QCD color} & \textbf{Standard $\C^3$ basis} &
\textbf{One-hot qubit state} & \textbf{Qubit basis index} \\
\midrule
Red   & $e_1 = (1,0,0)^T$ & $|100\rangle$ & 4 \\
Green & $e_2 = (0,1,0)^T$ & $|010\rangle$ & 2 \\
Blue  & $e_3 = (0,0,1)^T$ & $|001\rangle$ & 1 \\
\bottomrule
\end{tabular}
\caption{Basis mapping between standard QCD color space and one-hot three-qubit
states.  Qubit basis index uses binary encoding $4q_1 + 2q_2 + q_3$ (0-indexed).}
\label{tab:basis_mapping}
\end{table}

\subsection{Generator Mapping}

\begin{table}[h]
\centering
\begin{tabular}{ccc}
\toprule
\textbf{Index $a$} & \textbf{Standard generator $\frac{\lambda_a}{2}$} &
\textbf{Encoded generator $G_a = P\lambda_a P^\dagger$} \\
\midrule
$1$ & $r\leftrightarrow g$ off-diagonal, real    & $\Pcolor(\lambda_1)\Pcolor$, acts on $\{|100\rangle,|010\rangle\}$ \\
$2$ & $r\leftrightarrow g$ off-diagonal, imaginary & idem, imaginary \\
$3$ & $r$--$g$ color isospin diagonal             & $\Pcolor\,\mathrm{diag}(1,-1,0)\,\Pcolor$ \\
$4$ & $r\leftrightarrow b$ off-diagonal, real    & $\Pcolor(\lambda_4)\Pcolor$, acts on $\{|100\rangle,|001\rangle\}$ \\
$5$ & $r\leftrightarrow b$ off-diagonal, imaginary & idem, imaginary \\
$6$ & $g\leftrightarrow b$ off-diagonal, real    & $\Pcolor(\lambda_6)\Pcolor$, acts on $\{|010\rangle,|001\rangle\}$ \\
$7$ & $g\leftrightarrow b$ off-diagonal, imaginary & idem, imaginary \\
$8$ & color hypercharge diagonal                  & $\Pcolor\,\mathrm{diag}(1,1,-2)/\sqrt{3}\,\Pcolor$ \\
\bottomrule
\end{tabular}
\caption{Generator mapping.  $\Pcolor(\lambda_a)\Pcolor$ denotes the operator
$P \lambda_a P^\dagger$ on $\C^8$, which acts as $\lambda_a$ on $\Hcolor$ and
vanishes on $\Hcolor^\perp$.}
\label{tab:generator_mapping}
\end{table}

%──────────────────────────────────────────────────────────────────────────────
\section{Proof of Commutation Relations}
\label{sec:su3qubit:commutators}
%──────────────────────────────────────────────────────────────────────────────

\begin{theorem}[Encoded generators satisfy $\mathfrak{su}(3)$ on $\Hcolor$]
\label{thm:commutators}
For all $a,b \in \{1,\ldots,8\}$,
\[
  [G_a, G_b] = 2i f_{abc}\, G_c \quad \text{on } \Hcolor,
\]
where $f_{abc}$ are the totally antisymmetric $\mathfrak{su}(3)$ structure
constants.
\end{theorem}

\begin{proof}
Let $|v\rangle \in \Hcolor$ be arbitrary.  Write $|v\rangle = P\mathbf{v}$ for
a unique $\mathbf{v} \in \C^3$.  Since $P^\dagger P = I_3$ (isometry),
\begin{align*}
  G_a G_b |v\rangle
  &= (P\lambda_a P^\dagger)(P\lambda_b P^\dagger) P\mathbf{v} \\
  &= P\lambda_a (P^\dagger P)\lambda_b \mathbf{v} \\
  &= P\lambda_a \lambda_b \mathbf{v}.
\end{align*}
Therefore
\[
  [G_a, G_b]|v\rangle
  = P[\lambda_a, \lambda_b]\mathbf{v}
  = P (2i f_{abc}\lambda_c)\mathbf{v}
  = 2i f_{abc}\, G_c |v\rangle,
\]
using the standard $\mathfrak{su}(3)$ relation $[\lambda_a,\lambda_b] =
2i f_{abc}\lambda_c$ in the last step.

For $|v\rangle \in \Hcolor^\perp$: $G_a|v\rangle = P\lambda_a P^\dagger|v\rangle
= 0$ (since $P^\dagger|v\rangle = 0$), so $[G_a,G_b]|v\rangle = 0 =
2i f_{abc}G_c|v\rangle$.  The relation holds on all of $\Hthreeq$.
\end{proof}

\begin{remark}
The commutation relations hold on the \emph{entire} $\C^8$ (both
$\Hcolor$ and its complement), but the action is trivial on $\Hcolor^\perp$.
The algebra is realised faithfully on $\Hcolor \cong \C^3$.
\end{remark}

%──────────────────────────────────────────────────────────────────────────────
\section{Gauge Transformation Equivalence}
\label{sec:su3qubit:gauge_transform}
%──────────────────────────────────────────────────────────────────────────────

\begin{theorem}[Gauge transformations are compatible]
\label{thm:gauge_transform}
Let $\alpha^a(x)$ be real gauge parameters.  Define
\[
  U(x) = \exp\!\bigl(i\alpha^a(x)\tfrac{\lambda_a}{2}\bigr) \in SU(3),
  \qquad
  U_{3q}(x) = \exp\!\bigl(i\alpha^a(x)\tfrac{G_a}{2}\bigr) \in U(8).
\]
Then:
\begin{enumerate}
  \item $U_{3q}(x)\big|_{\Hcolor}$ is unitarily equivalent to $U(x)$ via $P$:
        \[
          U_{3q}(x)\, P = P\, U(x).
        \]
  \item $U_{3q}(x)\big|_{\Hcolor^\perp} = \mathrm{Id}|_{\Hcolor^\perp}$.
\end{enumerate}
\end{theorem}

\begin{proof}
For part (1): Since $G_a P = P \lambda_a$ (from Definition~\ref{def:Ga}) and
$\lambda_a$ matrices commute with this embedding, the exponential satisfies
\[
  \exp\!\bigl(i\alpha^a \tfrac{G_a}{2}\bigr) P
  = P \exp\!\bigl(i\alpha^a \tfrac{\lambda_a}{2}\bigr),
\]
which follows by expanding the matrix exponential and using $G_a P = P\lambda_a$
inductively (each power $(G_a)^n P = P(\lambda_a)^n$, verified by applying
$P^\dagger P = I_3$ at each step).

For part (2): $G_a|_{\Hcolor^\perp} = 0$, so
$\exp(i\alpha^a G_a/2)|_{\Hcolor^\perp} = e^0 = \mathrm{Id}$.
\end{proof}

%──────────────────────────────────────────────────────────────────────────────
\section{Covariant Derivative Equivalence}
\label{sec:su3qubit:cov_deriv}
%──────────────────────────────────────────────────────────────────────────────

\begin{theorem}[Covariant derivatives are equivalent on $\Hcolor$]
\label{thm:cov_deriv}
Let $\psi(x) \in \Hcolor$ be a color state (a section of the color bundle).
Write $\psi = P\varphi$ with $\varphi \in \C^3$.  Then
\[
  P^\dagger\bigl(D_\mu^{3q}\,\psi\bigr)
  = D_\mu\,\varphi,
\]
where $D_\mu = \partial_\mu + ig A_\mu^a \frac{\lambda_a}{2}$ is the standard
QCD covariant derivative on $\C^3$.
\end{theorem}

\begin{proof}
\begin{align*}
  D_\mu^{3q}\,\psi
  &= \partial_\mu(P\varphi) + ig A_\mu^a \frac{G_a}{2} (P\varphi) \\
  &= P(\partial_\mu\varphi) + ig A_\mu^a P\frac{\lambda_a}{2} \varphi \\
  &= P\!\left(\partial_\mu\varphi + ig A_\mu^a \frac{\lambda_a}{2} \varphi\right)
   = P(D_\mu\varphi).
\end{align*}
Applying $P^\dagger$ to both sides and using $P^\dagger P = I_3$ gives the
claim.
\end{proof}

\begin{corollary}[Gauge transformation of $D_\mu^{3q}$]
Under a local gauge transformation $U_{3q}(x)$, the encoded covariant
derivative transforms as
\[
  D_\mu^{3q} \;\longrightarrow\;
  U_{3q} D_\mu^{3q} U_{3q}^\dagger
  \quad\text{on } \Hcolor,
\]
consistent with the standard transformation law $D_\mu \to U D_\mu U^\dagger$
via the equivalence of Theorem~\ref{thm:cov_deriv}.
\end{corollary}

%──────────────────────────────────────────────────────────────────────────────
\section{Main Equivalence Theorem}
\label{sec:su3qubit:main_theorem}
%──────────────────────────────────────────────────────────────────────────────

\begin{tcolorbox}[colback=green!5!white,colframe=green!60!black,
  title=Main Theorem: SU(3) Gauge–Qubit Equivalence,fonttitle=\bfseries]
\begin{theorem}[SU(3) gauge--three-qubit equivalence]
\label{thm:main_equivalence}
Let $\Hcolor = \mathrm{span}\{|100\rangle,|010\rangle,|001\rangle\} \subset
\C^2\!\otimes\!\C^2\!\otimes\!\C^2$.  Then:

\begin{enumerate}
  \item \textbf{Generator equivalence.}
        The map $\lambda_a \mapsto G_a = P\lambda_a P^\dagger$ is a faithful
        representation of $\mathfrak{su}(3)$ on $\Hcolor$, i.e.\
        $[G_a, G_b] = 2i f_{abc}G_c$ on $\Hcolor$.

  \item \textbf{Gauge field equivalence.}
        The encoded gauge connection $A_\mu^{3q} = A_\mu^a G_a/2$ restricted to
        $\Hcolor$ is isomorphic to the standard connection $A_\mu = A_\mu^a
        \lambda_a/2$ via the isometry $P$.

  \item \textbf{Covariant derivative equivalence.}
        The encoded derivative $D_\mu^{3q}$ restricted to $\Hcolor$ is
        unitarily equivalent to the standard QCD covariant derivative $D_\mu$:
        \[
          P^\dagger D_\mu^{3q} P = D_\mu.
        \]

  \item \textbf{Gauge transformation equivalence.}
        Local SU(3) gauge transformations act consistently on both descriptions:
        $U_{3q}(x)\,P = P\,U(x)$.
\end{enumerate}

The equivalence is representation-level: it identifies the two descriptions as
the same $\mathfrak{su}(3)$ gauge structure realized in different but isometric
carrier spaces.
\end{theorem}
\end{tcolorbox}

%──────────────────────────────────────────────────────────────────────────────
\section{Leakage Sector Analysis}
\label{sec:su3qubit:leakage}
%──────────────────────────────────────────────────────────────────────────────

\subsection{Structure of the Leakage Sector}

The orthogonal complement of $\Hcolor$ in $\Hthreeq$ is the five-dimensional
space
\[
  \Hcolor^\perp = \mathrm{span}\{|000\rangle,\, |011\rangle,\, |101\rangle,\,
  |110\rangle,\, |111\rangle\}.
\]
Its physical content is:
\begin{itemize}
  \item $|000\rangle$: vacuum (no color excitation);
  \item $|011\rangle, |101\rangle, |110\rangle$: two-color (di-quark)
        excitations;
  \item $|111\rangle$: three-color (fully saturated) excitation.
\end{itemize}

\subsection{Gauge Action on the Leakage Sector}

All encoded generators $G_a$ vanish on $\Hcolor^\perp$
(Definition~\ref{def:Ga}), so gauge transformations act as the identity there.
The leakage sector is gauge-invariant but physically distinct from the color
subspace.

\subsection{Interpretation and Constraint Options}

Three interpretations are available; they should not be conflated:

\begin{enumerate}
  \item \textbf{Hard projection (physical constraint):} Impose $\Pcolor|\psi\rangle
        = |\psi\rangle$ as a superselection rule.  This recovers standard QCD
        color quantum numbers; the leakage sector is unphysical.

  \item \textbf{Auxiliary UBT degrees of freedom:} In the UBT context, the five
        leakage states may be reinterpreted as auxiliary biquaternionic degrees
        of freedom — e.g.\ the vacuum and multi-excitation modes of the $\Theta$
        field beyond the fundamental color triplet.  This interpretation is
        \textbf{speculative} and requires additional derivation from the UBT
        action; it is not established here.

  \item \textbf{Quantum simulation context:} When using the three-qubit space as
        a quantum-simulation register for SU(3) dynamics, the leakage sector
        must be actively suppressed (e.g.\ by energy penalty terms or
        error-correcting constraints) to maintain faithful simulation.
\end{enumerate}

\begin{tcolorbox}[colback=orange!10!white,colframe=orange!75!black,
  title=Warning: limits of the equivalence]
The equivalence of Theorem~\ref{thm:main_equivalence} holds
\textbf{strictly on $\Hcolor$}.

\medskip\noindent
Do \textbf{not} claim:
\begin{itemize}
  \item that the full Hilbert space $\Hthreeq = \C^8$ is identical to SU(3) or
        to any physical color Hilbert space;
  \item that the three-qubit construction provides a physical derivation of
        QCD until dynamics (Yang-Mills action), confinement, anomaly
        cancellation, and fermion representations are addressed;
  \item that the leakage sector $\Hcolor^\perp$ has a canonical UBT
        interpretation without further derivation.
\end{itemize}
\end{tcolorbox}

%──────────────────────────────────────────────────────────────────────────────
\section{Relation to Canonical UBT SU(3) Derivation}
\label{sec:su3qubit:relation_canonical}
%──────────────────────────────────────────────────────────────────────────────

The canonical UBT derivation of $\mathrm{SU}(3)_c$ proceeds via the involution
structure of $\C\otimes\HH$:
\[
  V_c = \ker(P_2 + \mathbf{1}) \cong \C^3, \qquad
  \mathrm{SU}(3) = \mathrm{Aut}(V_c, \langle\cdot,\cdot\rangle),
\]
(Theorems G.A--G.D; \texttt{canonical/su3\_derivation/}).  The present encoding
identifies $\Hcolor \cong \C^3$ via a different route (one-hot embedding in
$\Hthreeq$), but arrives at the same $\mathfrak{su}(3)$ algebra acting on the
same three-dimensional complex vector space.

The two constructions are therefore \textbf{equivalent representations of the
same Lie algebra}, realized via isometric embeddings of $\C^3$ into (different)
ambient eight-dimensional spaces:
\begin{align*}
  V_c = \ker(P_2+\mathbf{1}) \subset \C\otimes\HH \cong \C^8, \\
  \Hcolor = \mathrm{range}(P) \subset \C^2\otimes\C^2\otimes\C^2 \cong \C^8.
\end{align*}
Both identify the same group $\mathrm{SU}(3)$; neither replaces the canonical
derivation.  For the algebraic bridge between the two ambient $\C^8$ spaces see
the bridge document
\texttt{canonical/bridges/su3\_gauge\_qubit\_equivalence.tex}.

% End of su3_qubit_encoding.tex

\UBTSUThreeEndDocument
